\documentclass[11pt,a4paper]{article}
\usepackage{fontspec,amsmath,amssymb,amsthm,microtype}
\setmainfont{lmroman10-regular.otf}[BoldFont=lmroman10-bold.otf,ItalicFont=lmroman10-italic.otf,BoldItalicFont=lmroman10-bolditalic.otf]
\usepackage{xcolor}
\usepackage[most]{tcolorbox}
\usepackage{mathrsfs}
\providecommand{\lt}{<}
\providecommand{\gt}{>}
\setmonofont{DejaVuSansMono.ttf}[BoldFont=DejaVuSansMono-Bold.ttf,ItalicFont=DejaVuSansMono-Oblique.ttf,BoldItalicFont=DejaVuSansMono-BoldOblique.ttf,Scale=0.86]
\newfontfamily\leanfallback{STIXTwoMath-Regular.otf}[Scale=0.86]
\definecolor{leanrule}{gray}{0.62}
\definecolor{leanback}{gray}{0.962}
\definecolor{leantitle}{gray}{0.22}
\definecolor{leanmuted}{gray}{0.38}
\newtcolorbox{leanbox}[1]{enhanced,breakable,colback=leanback,colframe=leanrule,
  boxrule=0.4pt,arc=0.6mm,left=2.4mm,right=2.4mm,top=1.8mm,bottom=1.8mm,
  before skip=0.8em,after skip=1em,fonttitle=\small\bfseries,coltitle=leantitle,
  colbacktitle=leanback,titlerule=0.3pt,title={#1}}
\newlength{\leancharwidth}
\newenvironment{leancode}{\par\smallskip\begingroup\ttfamily\footnotesize\raggedright
  \settowidth{\leancharwidth}{M}\setlength{\parskip}{0pt}\setlength{\parindent}{0pt}
  \frenchspacing}{\par\endgroup\smallskip}
\newcommand{\leanline}[2]{\hangindent=\dimexpr#1\leancharwidth+4\leancharwidth\relax
  \hangafter=1\noindent\hspace*{#1\leancharwidth}#2\par}
\newcommand{\leansource}[1]{\par{\footnotesize\color{leanmuted}#1}\par}
\newcommand{\leanlabel}[1]{\par\smallskip\noindent{\small\bfseries #1}\par}
\usepackage[margin=25mm]{geometry}
\usepackage[unicode,hidelinks]{hyperref}
\hypersetup{pdfauthor={LangevinLean Project}}
\setlength{\emergencystretch}{3em}
\setlength{\parskip}{0.4em}
\setlength{\parindent}{0pt}
\allowdisplaybreaks[2]
\raggedbottom
\title{\texorpdfstring{A Cholesky factor for the Gaussian response covariance}{A Cholesky factor for the Gaussian response covariance}}
\author{LangevinLean Project}
\date{}
\begin{document}\maketitle
Draft mathematical exposition; these notes may contain mistakes.


A Gaussian law specified through its covariance must agree with the specified Gaussian construction to multiply independent standard normals by a Cholesky factor. This section proves that exact identity, constructs the positive-diagonal lower-triangular \(2\times2\) factor, and applies it to the strictly positive exponential Gram covariance.



An explicit lower-triangular factor of the response Gram covariance maps independent scalar standard normals to the exact sampler innovation law. The proof includes positive pivots and Gaussian measure equality.


\section{\texorpdfstring{Identify every covariance factor by mean and covariance}{Identify every covariance factor by mean and covariance}}


If \(AA^\mathsf T=S\), then applying \(A\) to a standard Gaussian gives a centered Gaussian with covariance \(S\). View the matrix as a linear operator on Euclidean space. Its transpose is the adjoint, and matrix multiplication is composition. The mean passes through the continuous linear map and remains zero. The covariance pushforward is \(\langle A^\mathsf Tx,A^\mathsf Ty\rangle=\langle x,AA^\mathsf Ty\rangle\), which is the target covariance. Uniqueness of Gaussian measures with equal mean and covariance proves equality of measures, not only equality of moments.



A deterministic linear image of independent standard normals is jointly Gaussian. Its covariance is the factor times its transpose. Equality of centered Gaussian covariances determines equality of measures, which is why any verified covariance factor gives the exact response law.

\begin{leanbox}{Lean: \texttt{gaussian\_\allowbreak{}factor\_\allowbreak{}hasLaw}}
\leansource{\texttt{gaussian\_\allowbreak{}factor\_\allowbreak{}hasLaw} \quad KineticLangevin/\allowbreak{}GaussianFactor.lean, lines 12–32}
\begin{leancode}
\leanline{0}{theorem\ gaussian\_factor\_hasLaw\ (S\ A\ :\ Matrix\ ι\ ι\ ℝ)\ (hS\ :\ S.PosSemidef)\ (hA\ :\ A*Aᴴ=S)\ :}
\leanline{4}{HasLaw\ (toEuclideanCLM\ ({\leanfallback 𝕜}:=ℝ)\ A)\ (multivariateGaussian\ 0\ S)}
\leanline{6}{(stdGaussian\ (EuclideanSpace\ ℝ\ ι))}
\end{leancode}
\leanlabel{How the Lean proof follows the argument}
\texttt{gaussian\_\allowbreak{}factor\_\allowbreak{}has\allowbreak{}Law} gives the formal counterpart of this step. The main intermediate facts are \texttt{L} in that order. Named proof applications include \texttt{Is\allowbreak{}Gaussian.\allowbreak{}ext}, \texttt{hi}, \texttt{congr\allowbreak{}Arg}; the adjacent proof shows their hypotheses and order of use.
\end{leanbox}

\section{\texorpdfstring{Construct and verify the positive-diagonal Cholesky factor}{Construct and verify the positive-diagonal Cholesky factor}}


For a positive definite symmetric \(2\times2\) matrix with entries \(a=S_{00}\), \(b=S_{10}\), \(d=S_{11}\), define \(L=\begin{pmatrix}\sqrt a&0\\b/\sqrt a&\sqrt{d-b^2/a}\end{pmatrix}\). Positive definiteness gives \(a>0\) and \(ad-b^2>0\), hence a strictly positive Schur complement. The factor is lower triangular with positive diagonal. Expanding its four matrix entries, using the two square-root-square identities and symmetry, gives \(LL^\mathsf T=S\). The general Gaussian-factor theorem therefore proves that this triangular factor generates exactly the target Gaussian measure.



The first pivot is positive by positive definiteness. The determinant condition makes the Schur complement positive, so both diagonal square roots are real and nonzero. Expanding the lower-triangular product verifies each matrix entry without an implicit numerical factorization assumption.

\begin{leanbox}{Lean: \texttt{choleskyTwo}, \texttt{choleskyTwo\_\allowbreak{}schur\_\allowbreak{}pos}, \texttt{choleskyTwo\_\allowbreak{}triangular} and 2 more}
\leansource{\texttt{choleskyTwo} \quad KineticLangevin/\allowbreak{}GaussianFactor.lean, lines 35–36}
\begin{leancode}
\leanline{0}{def\ choleskyTwo\ (S\ :\ Matrix\ (Fin\ 2)\ (Fin\ 2)\ ℝ)\ :\ Matrix\ (Fin\ 2)\ (Fin\ 2)\ ℝ\ :=}
\leanline{2}{!![Real.sqrt\ (S\ 0\ 0),0;S\ 1\ 0/Real.sqrt\ (S\ 0\ 0),Real.sqrt\ (S\ 1\ 1-(S\ 1\ 0)\textasciicircum{}2/S\ 0\ 0)]}
\end{leancode}
\leansource{\texttt{choleskyTwo\_\allowbreak{}schur\_\allowbreak{}pos} \quad KineticLangevin/\allowbreak{}GaussianFactor.lean, lines 39–47}
\begin{leancode}
\leanline{0}{theorem\ choleskyTwo\_schur\_pos\ (S\ :\ Matrix\ (Fin\ 2)\ (Fin\ 2)\ ℝ)\ (hS\ :\ S.PosDef)\ :}
\leanline{4}{0<S\ 1\ 1-(S\ 1\ 0)\textasciicircum{}2/S\ 0\ 0}
\end{leancode}
\leansource{\texttt{choleskyTwo\_\allowbreak{}triangular} \quad KineticLangevin/\allowbreak{}GaussianFactor.lean, lines 50–59}
\begin{leancode}
\leanline{0}{theorem\ choleskyTwo\_triangular\ (S\ :\ Matrix\ (Fin\ 2)\ (Fin\ 2)\ ℝ)\ (hS\ :\ S.PosDef)\ :}
\leanline{4}{(choleskyTwo\ S).IsLowerTriangular\ ∧\ ∀\ i,0<choleskyTwo\ S\ i\ i}
\end{leancode}
\leansource{\texttt{choleskyTwo\_\allowbreak{}factor} \quad KineticLangevin/\allowbreak{}GaussianFactor.lean, lines 62–78}
\begin{leancode}
\leanline{0}{theorem\ choleskyTwo\_factor\ (S\ :\ Matrix\ (Fin\ 2)\ (Fin\ 2)\ ℝ)\ (hS\ :\ S.PosDef)\ :}
\leanline{4}{choleskyTwo\ S*(choleskyTwo\ S)ᴴ=S}
\end{leancode}
\leansource{\texttt{choleskyTwo\_\allowbreak{}hasLaw} \quad KineticLangevin/\allowbreak{}GaussianFactor.lean, lines 81–84}
\begin{leancode}
\leanline{0}{theorem\ choleskyTwo\_hasLaw\ (S\ :\ Matrix\ (Fin\ 2)\ (Fin\ 2)\ ℝ)\ (hS\ :\ S.PosDef)\ :}
\leanline{4}{HasLaw\ (toEuclideanCLM\ ({\leanfallback 𝕜}:=ℝ)\ (choleskyTwo\ S))\ (multivariateGaussian\ 0\ S)}
\leanline{6}{(stdGaussian\ (EuclideanSpace\ ℝ\ (Fin\ 2)))}
\end{leancode}
\leanlabel{How the Lean proof follows the argument}
\texttt{cholesky\allowbreak{}Two} names the mathematical object defined above. Its defining expression appears in the adjacent statement; subsequent identities may unfold this definition directly.




\texttt{cholesky\allowbreak{}Two\_\allowbreak{}schur\_\allowbreak{}pos} gives the formal counterpart of this step. The main intermediate facts are \texttt{ha}, \texttt{hd}, \texttt{hsym} in that order. Its remaining equalities are discharged by the simplifications displayed in the proof.




\texttt{cholesky\allowbreak{}Two\_\allowbreak{}triangular} gives the formal counterpart of this step.




\texttt{cholesky\allowbreak{}Two\_\allowbreak{}factor} gives the formal counterpart of this step. The main intermediate facts are \texttt{ha}, \texttt{hs}, \texttt{hq}, \texttt{hsq}, \texttt{hsym}, \texttt{he} in that order. Named proof applications include \texttt{hqeq}; the adjacent proof shows their hypotheses and order of use.




\texttt{cholesky\allowbreak{}Two\_\allowbreak{}has\allowbreak{}Law} gives the formal counterpart of this step.
\end{leanbox}

\section{\texorpdfstring{Generate the exponential innovation from independent normals}{Generate the exponential innovation from independent normals}}


The two response kernels \(\psi_0,\psi_1\) have a strictly positive Gram covariance on every positive-length interval. Apply the explicit Cholesky construction to that matrix. The resulting output law is precisely the two-response Gaussian law, already identified with the Brownian convolution pair used by the sampler. Finally, the product of two real standard Gaussian measures pushes forward to the Euclidean standard Gaussian under the coordinate identification. Composing the two exact laws proves the paper’s realization starting from two independent scalar normal draws. This proves the covariance-factor identity for the shared exponential integrator and its particle specialization.



The response Gram matrix is strictly positive for positive step length. Its explicit triangular factor therefore transforms two independent scalar standard normals into the velocity–position innovation pair. This gives an exact realization of the same law used by the Brownian coupling proof.

\begin{leanbox}{Lean: \texttt{klmc\_\allowbreak{}cholesky\_\allowbreak{}hasLaw}, \texttt{klmc\_\allowbreak{}cholesky\_\allowbreak{}independent\_\allowbreak{}normals}}
\leansource{\texttt{klmc\_\allowbreak{}cholesky\_\allowbreak{}hasLaw} \quad KineticLangevin/\allowbreak{}GaussianFactor.lean, lines 87–90}
\begin{leancode}
\leanline{0}{theorem\ klmc\_cholesky\_hasLaw\ (γ\ h\ :\ ℝ)\ (hh\ :\ 0<h)\ :}
\leanline{4}{HasLaw\ (toEuclideanCLM\ ({\leanfallback 𝕜}:=ℝ)\ (choleskyTwo\ (klmcCovariance\ γ\ h)))\ (klmcGaussian\ γ\ h)}
\leanline{6}{(stdGaussian\ (EuclideanSpace\ ℝ\ (Fin\ 2)))}
\end{leancode}
\leansource{\texttt{klmc\_\allowbreak{}cholesky\_\allowbreak{}independent\_\allowbreak{}normals} \quad KineticLangevin/\allowbreak{}GaussianFactor.lean, lines 93–99}
\begin{leancode}
\leanline{0}{theorem\ klmc\_cholesky\_independent\_normals\ (γ\ h\ :\ ℝ)\ (hh\ :\ 0<h)\ :}
\leanline{4}{HasLaw\ (fun\ z\ :\ Fin\ 2→ℝ\ =>\ toEuclideanCLM\ ({\leanfallback 𝕜}:=ℝ)\ (choleskyTwo\ (klmcCovariance\ γ\ h))\ (WithLp.toLp\ 2\ z))}
\leanline{6}{(klmcGaussian\ γ\ h)\ (Measure.pi\ (fun\ \_\ :\ Fin\ 2\ =>\ gaussianReal\ 0\ 1))}
\end{leancode}
\leanlabel{How the Lean proof follows the argument}
\texttt{klmc\_\allowbreak{}cholesky\_\allowbreak{}has\allowbreak{}Law} gives the formal counterpart of this step.




\texttt{klmc\_\allowbreak{}cholesky\_\allowbreak{}independent\_\allowbreak{}normals} gives the formal counterpart of this step. The main intermediate facts are \texttt{hl} in that order.
\end{leanbox}

\end{document}
